\documentclass[12pt]{article}
\usepackage{amsmath,amssymb}

\begin{document}
\section*{{2024 AIME I Problems/Problem 9}}
Source: AoPS Wiki

\subsection*{Solution}
A quadrilateral is a rhombus if and only if its two diagonals bisect each other and are perpendicular to each other. The first condition is automatically satisfied because of the hyperbola's symmetry about the origin. To satisfy the second condition, we set $BD$ as the line $y = mx$ and $AC$ as $y = -\frac{1}{m}x.$ Because the hyperbola has asymptotes of slopes $\pm \frac{\sqrt6}{\sqrt5},$ we have $m, -\frac{1}{m} \in \left(-\frac{\sqrt6}{\sqrt5}, \frac{\sqrt6}{\sqrt5}\right).$ This gives us $m^2 \in \left(\frac{5}{6}, \frac{6}{5}\right).$ Plugging $y = mx$ into the equation for the hyperbola yields $x^2 = \frac{120}{6-5m^2}$ and $y^2 = \frac{120m^2}{6-5m^2}.$ By symmetry of the hyperbola, we know that $\left(\frac{BD}{2}\right)^2 = x^2 + y^2,$ so we wish to find a lower bound for $x^2 + y^2 = 120\left(\frac{1+m^2}{6-5m^2}\right).$ This is equivalent to minimizing $\frac{1+m^2}{6-5m^2} = -\frac{1}{5} + \frac{11}{5(6-5m^2)}$ . It's then easy to see that this expression increases with $m^2,$ so we plug in $m^2 = \frac{5}{6}$ to get $x^2+y^2 > 120,$ giving $BD^2 > \boxed{480}.$

\end{document}
